\subsection{一次変換}

ベクトル空間 $\mathbb{R}^n$ からベクトル空間 $\mathbb{R}^n$ への写像を
考える。写像 $f\colon\mathbb{R}^n\to\mathbb{R}^n$ が次の二つの性質を
満たすとき、$f$ を\textbf{一次変換}（線形写像）という。
\begin{itemize}
 \item 加法をしてから変換しても、変換してから加法をしても同じである：
 \begin{equation}
  f(\vect{v}+\vect{w})=f(\vect{v})+f(\vect{w}).
 \end{equation}
 \item スカラー倍してから変換しても、変換してからスカラー倍しても同じである：
 \begin{equation}
  f(a\vect{v})=a f(\vect{v}).
 \end{equation}
\end{itemize}

\begin{example}
変換
\[
 A=\begin{pmatrix}1&-1\\2&1\end{pmatrix},\qquad
 \vect{v}=\begin{pmatrix}2\\1\end{pmatrix},\qquad
 \vect{w}=\begin{pmatrix}-1\\1\end{pmatrix}
\]
について、加法とスカラー倍が保たれることを確かめる。
\begin{align*}
 A(\vect{v}+\vect{w})
 &=\begin{pmatrix}1&-1\\2&1\end{pmatrix}
   \left\{\begin{pmatrix}2\\1\end{pmatrix}
   +\begin{pmatrix}-1\\1\end{pmatrix}\right\}
 =\begin{pmatrix}1&-1\\2&1\end{pmatrix}
  \begin{pmatrix}1\\2\end{pmatrix}
 =\begin{pmatrix}-1\\4\end{pmatrix},\\
 A\vect{v}+A\vect{w}
 &=\begin{pmatrix}1&-1\\2&1\end{pmatrix}\begin{pmatrix}2\\1\end{pmatrix}
  +\begin{pmatrix}1&-1\\2&1\end{pmatrix}\begin{pmatrix}-1\\1\end{pmatrix}\\
 &=\begin{pmatrix}1\\5\end{pmatrix}+\begin{pmatrix}-2\\-1\end{pmatrix}
 =\begin{pmatrix}-1\\4\end{pmatrix}.
\end{align*}
また、
\begin{align*}
 A(2\vect{v})
 &=\begin{pmatrix}1&-1\\2&1\end{pmatrix}\begin{pmatrix}4\\2\end{pmatrix}
 =\begin{pmatrix}2\\10\end{pmatrix},\\
 2A\vect{v}
 &=2\begin{pmatrix}1&-1\\2&1\end{pmatrix}\begin{pmatrix}2\\1\end{pmatrix}
 =2\begin{pmatrix}1\\5\end{pmatrix}
 =\begin{pmatrix}2\\10\end{pmatrix}.
\end{align*}
\end{example}

\subsubsection{直交変換}

内積を保存する一次変換を\textbf{直交変換}という。すなわち、
\begin{align}
 (A\vect{v})\mathbin{\boldsymbol{\cdot}}(A\vect{w})
 &=(A\vect{v})^{\mathsf T}(A\vect{w})\notag\\
 &=\vect{v}^{\mathsf T}A^{\mathsf T}A\vect{w}
 =\vect{v}^{\mathsf T}\vect{w}
\end{align}
を満たす変換である。したがって、直交変換の表現行列 $A$ は
\begin{equation}
 \boxed{A^{\mathsf T}A=I}
\end{equation}
を満たす。このような行列を\textbf{直交行列}という。

直交変換について、次が成り立つ。
\begin{itemize}
 \item 逆行列は転置行列である：$A^{-1}=A^{\mathsf T}$。
 \item 行列式は $\det A=\pm1$ である。実際、
 \[
  1=\det I=\det(A^{\mathsf T}A)
   =(\det A^{\mathsf T})(\det A)=(\det A)^2.
 \]
\end{itemize}

\begin{example}
\begin{equation}
 A=\begin{pmatrix}
 \dfrac1{\sqrt2}&-\dfrac1{\sqrt2}\\[2mm]
 \dfrac1{\sqrt2}& \dfrac1{\sqrt2}
 \end{pmatrix}
\end{equation}
とする。このとき
\[
 \det A=\frac12-\left(-\frac12\right)=1,
\]
また
\[
 A^{\mathsf T}A
 =\begin{pmatrix}
 \dfrac1{\sqrt2}&\dfrac1{\sqrt2}\\[2mm]
 -\dfrac1{\sqrt2}&\dfrac1{\sqrt2}
 \end{pmatrix}
 \begin{pmatrix}
 \dfrac1{\sqrt2}&-\dfrac1{\sqrt2}\\[2mm]
 \dfrac1{\sqrt2}&\dfrac1{\sqrt2}
 \end{pmatrix}
 =\begin{pmatrix}1&0\\0&1\end{pmatrix}
\]
であるから、$A$ は直交変換である。

さらに
\[
 \vect{v}=\begin{pmatrix}2\\1\end{pmatrix},\qquad
 \vect{w}=\begin{pmatrix}-1\\1\end{pmatrix}
\]
とすると、
\[
 \vect{v}\mathbin{\boldsymbol{\cdot}}\vect{w}
 =\vect{v}^{\mathsf T}\vect{w}=-1,
\]
一方、
\[
 A\vect{v}=\begin{pmatrix}1/\sqrt2\\3/\sqrt2\end{pmatrix},\qquad
 A\vect{w}=\begin{pmatrix}-2/\sqrt2\\0\end{pmatrix}
\]
なので、
\[
 (A\vect{v})\mathbin{\boldsymbol{\cdot}}(A\vect{w})=-1
\]
となり、内積が保存されている。
\end{example}

この $A$ はどのような変換であろうか。標準基底の像は
\[
 \unit{x}'=A\unit{x}
 =\begin{pmatrix}1/\sqrt2\\1/\sqrt2\end{pmatrix},\qquad
 \unit{y}'=A\unit{y}
 =\begin{pmatrix}-1/\sqrt2\\1/\sqrt2\end{pmatrix}
\]
である。

\begin{figure}[htbp]
\centering
\begin{tikzpicture}[>=Stealth,scale=2.2]
 \draw[->] (-1.25,0)--(1.45,0) node[right] {$x$};
 \draw[->] (0,-.25)--(0,1.45) node[above] {$y$};
 \draw[densely dashed,gray] (0,0) circle (1);
 \draw[->,thick] (0,0)--(1,0) node[below right] {$\unit{x}$};
 \draw[->,thick] (0,0)--(0,1) node[above right] {$\unit{y}$};
 \draw[->,very thick,blue] (0,0)--({sqrt(2)/2},{sqrt(2)/2})
   node[right] {$\unit{x}'$};
 \draw[->,very thick,red] (0,0)--({-sqrt(2)/2},{sqrt(2)/2})
   node[above left] {$\unit{y}'$};
 \draw[->] (.38,0) arc[start angle=0,end angle=45,radius=.38]
   node[midway,right] {$\frac{\pi}{4}$};
 \draw[->] (0,.48) arc[start angle=90,end angle=135,radius=.48]
   node[midway,above] {$\frac{\pi}{4}$};
 \node[below left] at (0,0) {$O$};
\end{tikzpicture}
\caption{$\unit{x}$ と $\unit{y}$ をそれぞれ原点のまわりに
$\pi/4$ だけ回転する変換}
\end{figure}

\subsubsection{2次元の回転}

標準基底
\[
 \unit{x}=\begin{pmatrix}1\\0\end{pmatrix},\qquad
 \unit{y}=\begin{pmatrix}0\\1\end{pmatrix}
\]
を、原点のまわりに反時計回りに角 $\theta$ だけ回転する。
\[
 \unit{x}'=\begin{pmatrix}\cos\theta\\\sin\theta\end{pmatrix},\qquad
 \unit{y}'=\begin{pmatrix}-\sin\theta\\\cos\theta\end{pmatrix}.
\]

\begin{figure}[htbp]
\centering
\begin{tikzpicture}[>=Stealth,scale=1.8]
 \def\ang{32}
 \draw[->] (-1.35,0)--(1.65,0) node[right] {$x$};
 \draw[->] (0,-.35)--(0,1.5) node[above] {$y$};
 \draw[->,thick] (0,0)--(1.15,0) node[below right] {$\unit{x}$};
 \draw[->,very thick,blue] (0,0)--(\ang:1.25) node[right] {$\unit{x}'$};
 \draw[->,very thick,red] (0,0)--({90+\ang}:1.25) node[above left] {$\unit{y}'$};
 \draw[->] (.46,0) arc[start angle=0,end angle=\ang,radius=.46]
   node[midway,right] {$\theta$};
 \draw[->] (0,.56) arc[start angle=90,end angle={90+\ang},radius=.56]
   node[midway,above] {$\theta$};
 \node[above right] at (0,.95) {$\unit{y}$};
\end{tikzpicture}
\caption{標準基底を反時計回りに角 $\theta$ だけ回転する}
\end{figure}

回転の一次変換を $R(\theta)$ とすると、
\[
 \begin{pmatrix}\cos\theta\\\sin\theta\end{pmatrix}
 =R(\theta)\begin{pmatrix}1\\0\end{pmatrix},\qquad
 \begin{pmatrix}-\sin\theta\\\cos\theta\end{pmatrix}
 =R(\theta)\begin{pmatrix}0\\1\end{pmatrix}.
\]
これらをまとめると、
\begin{equation}
 \boxed{R(\theta)=
 \begin{pmatrix}
  \cos\theta&-\sin\theta\\
  \sin\theta& \cos\theta
 \end{pmatrix}}
\end{equation}
となる。

次に、ベクトル $\vect{v}$ に角 $\alpha$ と角 $\beta$ の変換を順に
作用させると、
\[
 \vect{v}''=R(\beta)\vect{v}'=R(\beta)R(\alpha)\vect{v}.
\]
図から、これは $\vect{v}$ に角 $\alpha+\beta$ の変換を作用させる
ことと同じなので、
\begin{equation}
 \boxed{R(\alpha+\beta)=R(\beta)R(\alpha)}
\end{equation}
である。

\begin{figure}[htbp]
\centering
\begin{tikzpicture}[>=Stealth,scale=1.65]
 \draw[->] (-.25,0)--(2.25,0);
 \draw[->,thick] (0,0)--(18:1.65) node[right] {$\vect{v}$};
 \draw[->,very thick,blue] (0,0)--(50:1.75) node[right] {$\vect{v}'$};
 \draw[->,very thick,red] (0,0)--(80:1.85) node[above] {$\vect{v}''$};
 \draw[->] (18:.7) arc[start angle=18,end angle=50,radius=.7]
   node[midway,right] {$\alpha$};
 \draw[->] (50:1.0) arc[start angle=50,end angle=80,radius=1]
   node[midway,above right] {$\beta$};
\end{tikzpicture}
\caption{回転の合成}
\end{figure}

\paragraph{演習}
$R(\alpha+\beta)=R(\alpha)R(\beta)$ から、三角関数の加法定理
\begin{align*}
 \cos(\alpha+\beta)&=\cos\alpha\cos\beta-\sin\alpha\sin\beta,\\
 \sin(\alpha+\beta)&=\sin\alpha\cos\beta+\cos\alpha\sin\beta
\end{align*}
を求めよ。

\subsubsection{基底の変換}

たとえば、3次元の正規直交基底を
\[
 (\unit{1},\unit{2},\unit{3})
 \longrightarrow(\unit{1}',\unit{2}',\unit{3}')
\]
と変換する。このとき、ベクトル $\vect{v}$ の成分はどのように
変換されるだろうか。

ベクトルそのものは基底を変換しても変化しないので、
\begin{equation}
 \vect{v}
 =(\unit{1}\ \unit{2}\ \unit{3})
  \begin{pmatrix}x_1\\x_2\\x_3\end{pmatrix}
 =(\unit{1}'\ \unit{2}'\ \unit{3}')
  \begin{pmatrix}x_1'\\x_2'\\x_3'\end{pmatrix}
\end{equation}
である。ここで、基底の変換を $A$ とすると、
\[
 (\unit{1}'\ \unit{2}'\ \unit{3}')
 =(\unit{1}\ \unit{2}\ \unit{3})A.
\]
したがって
\[
 (\unit{1}'\ \unit{2}'\ \unit{3}')
 \begin{pmatrix}x_1'\\x_2'\\x_3'\end{pmatrix}
 =(\unit{1}\ \unit{2}\ \unit{3})AA^{-1}
 \begin{pmatrix}x_1\\x_2\\x_3\end{pmatrix}
\]
となり、新しい基底でのベクトルの成分は
\begin{equation}
 \boxed{
 \begin{pmatrix}x_1'\\x_2'\\x_3'\end{pmatrix}
 =A^{-1}\begin{pmatrix}x_1\\x_2\\x_3\end{pmatrix}}
\end{equation}
と変換される。

\begin{example}[座標軸の反転]
$y$ 軸を反転した座標系を考える。
\[
 (\unit{x}\ \unit{y})
 =\begin{pmatrix}1&0\\0&1\end{pmatrix}
 \longrightarrow
 \begin{pmatrix}1&0\\0&-1\end{pmatrix}
 =(\unit{x}'\ \unit{y}').
\]

\begin{figure}[htbp]
\centering
\begin{tikzpicture}[>=Stealth]
 \begin{scope}[xshift=-3cm]
  \draw[->] (0,0)--(1.6,0) node[right] {$x$};
  \draw[->] (0,0)--(0,1.4) node[above] {$y$};
 \end{scope}
 \draw[->,thick] (-.8,.7)--(.8,.7) node[midway,above] {$y$軸を反転};
 \begin{scope}[xshift=3cm]
  \draw[->] (0,0)--(1.6,0) node[right] {$x'$};
  \draw[->] (0,1.1)--(0,-1.4) node[below] {$y'$};
 \end{scope}
\end{tikzpicture}
\caption{$y$ 軸を反転した座標系}
\end{figure}

基底の変換行列は
\[
 A=\begin{pmatrix}1&0\\0&-1\end{pmatrix}
\]
である。したがって、位置ベクトルの成分は
\begin{equation}
 \begin{pmatrix}x'\\y'\end{pmatrix}
 =A^{-1}\begin{pmatrix}x\\y\end{pmatrix}
 =\begin{pmatrix}1&0\\0&-1\end{pmatrix}
  \begin{pmatrix}x\\y\end{pmatrix}
\end{equation}
と変換される。
\end{example}
